\documentclass[11pt,a4paper]{article}

\usepackage[T1]{fontenc}
\usepackage[utf8]{inputenc}
\usepackage[margin=1in]{geometry}
\usepackage{amsmath,amssymb}
\usepackage{booktabs}
\usepackage{array}
\usepackage{multirow}
\usepackage{graphicx}
\usepackage{url}
\usepackage{hyperref}
\usepackage[numbers,sort&compress]{natbib}
\usepackage{setspace}

\geometry{margin=1in}
\onehalfspacing

\hypersetup{
  colorlinks=true,
  linkcolor=black,
  citecolor=black,
  urlcolor=blue
}

\title{Fact Validator: A Transparent, Production-Ready Fact-Checking System with Auditable Credibility Scoring, Optional Debate Arbitration, and Cost-Aware Evidence Retrieval}
\author{Fact Validator Project Team}
\date{\today}

\begin{document}

\maketitle

\begin{abstract}
Fact Validator is an end-to-end fact-checking system that accepts a URL or free text, decomposes the input into claims, retrieves web evidence, scores source credibility with a human-auditable rubric, and returns verdicts with explanations and confidence estimates. The system is designed as a production artifact rather than a narrowly tuned benchmark model: it includes caching, feature flags, health checks, structured logging, and an optional LLM debate mode for uncertain cases. We evaluate the current implementation on a stratified 51-claim test split derived from a diversified benchmark and compare the full pipeline against simple baselines and ablations. On this split, the full system achieves 21.6\% accuracy and 0.212 macro-F1, while the strongest simple baseline reaches 37.3\% accuracy. The results therefore do not support claims of benchmark superiority. However, the credibility component contributes a measurable lift over its ablation, and the operational layer reduces estimated monthly SerpAPI cost by 71.4\% while preserving reproducibility and transparency. The main contribution of this work is a defensible system design for source-aware misinformation detection, not a claim that the current benchmark setup establishes state-of-the-art factuality performance.

\textbf{Keywords:} fact checking, misinformation detection, source credibility, explainable AI, LLM debate, evidence retrieval, production systems
\end{abstract}

\section{Introduction}

Automated fact checking is attractive because the volume of online claims exceeds what human reviewers can handle, yet most published systems still trade away transparency, operational robustness, or both. In practice, a fact-checking system must do more than classify claims: it must find evidence, decide what to trust, justify verdicts, remain stable when external services fail, and do so at acceptable cost. These requirements are especially important in deployment settings where users need to understand why a claim was marked supported, refuted, or unsupported.

Fact Validator was built to address this gap. The system combines retrieval-based verification, auditable source credibility scoring, optional LLM debate arbitration, caching, and production-grade error handling in a single pipeline. The design is inspired by work on hallucination detection, tool-augmented factuality, and faithful error correction \citep{galitsky2023truthometer,grave2020factool,kubota2023zero,thorne2018fever,ji2023hallucination,wen2024abstention}.

This manuscript makes two claims. First, Fact Validator is a transparent and deployable system for fact-checking source material at the claim level. Second, the current evaluation is useful for system characterization but not strong enough to support broad superiority claims over simple baselines. That distinction matters: the benchmark evidence is honest, but limited.

\section{Related Work}

Truth-O-Meter frames hallucination detection as evidence-backed reasoning over web sources and uses defeasible logic to resolve conflicts \citep{galitsky2023truthometer}. FacTool emphasizes modular tool augmentation for factuality verification and highlights the value of composing retrieval, external APIs, and inference tools \citep{grave2020factool}. Zero-shot faithful factual error correction shows that faithful correction can be framed as an external evidence task rather than a closed-world language modeling problem \citep{kubota2023zero}.

Fact Validator differs from these lines of work in three ways. First, the credibility layer is explicitly human-auditable rather than learned as a black box. Second, the system is packaged as a production application with persistent storage, observability, caching, and fallback behavior. Third, the evaluation is reported with operational metrics alongside model metrics, so the paper treats latency and cost as first-class outcomes rather than side effects.

The broader misinformation and hallucination literature also underscores the importance of abstention, calibration, and failure-aware evaluation \citep{ji2023hallucination,wen2024abstention}. Those concerns are central here because a system that sounds confident while being poorly calibrated is not suitable for practical fact checking.

\section{System Design}

Fact Validator follows a nine-stage pipeline:

\begin{enumerate}
  \item content extraction from a URL or acceptance of raw text,
  \item claim decomposition into fact-like statements,
  \item web evidence retrieval through SerpAPI,
  \item source credibility scoring,
  \item semantic reranking of retrieved evidence,
  \item baseline verdict classification,
  \item optional LLM debate arbitration,
  \item sentiment and bias analysis,
  \item final misinformation scoring and persistence.
\end{enumerate}

The architecture is intentionally conservative. The system does not rely on a single generative model to produce a verdict; instead, it combines retrieval, scoring, and rule-based aggregation so that each stage can be inspected independently.

\subsection{Credibility Scoring}

The credibility score is intentionally simple and falsifiable:

\[
\mathrm{Score}(d) = \mathrm{clip}_{[0,100]}\bigl(50 + b(d) + p(d)\bigr)
\]

where the baseline is neutral, $b(d)$ adds reputation bonuses for reliable domains, and $p(d)$ applies penalties for known low-trust hosts or platform-like sources. The exact rubric is intentionally explicit so that reviewers can argue over the mapping rather than reverse-engineer a learned trust model.

\subsection{Optional Debate Arbitration}

For uncertain cases, the system can invoke a Prover/Skeptic/Judge pattern through a local LLM endpoint. The Prover argues for support, the Skeptic surfaces weaknesses, and the Judge synthesizes the result into a final verdict. If the model is unavailable, the pipeline degrades gracefully to the baseline verifier.

\subsection{Cost-Aware Retrieval}

Evidence search is expensive and rate-limited, so the system memoizes normalized claims with a 24-hour TTL. This makes repeated analyses faster and reduces duplicate web queries. In the current configuration, the estimated monthly SerpAPI cost falls from USD 77.00 to USD 22.00 for an assumed 1000 claims, a 71.4\% reduction.

\section{Evaluation Setup}

\subsection{Benchmark}

The current benchmark comes from a diversified v2 dataset split into 143 training claims, 46 validation claims, and 51 test claims. The split is stratified by label and difficulty, but the benchmark still contains synthetic patterns and duplicates. For that reason, the results are appropriate for system validation and comparative analysis, but they should not be treated as a final scientific benchmark.

\subsection{Baselines and Ablations}

We compare the full pipeline against simple heuristics and component ablations:

\begin{itemize}
  \item random,
  \item majority,
  \item length heuristic,
  \item keyword heuristic,
  \item sentiment heuristic,
  \item ablate credibility,
  \item ablate semantic reranking,
  \item ablate debate,
  \item ablate quality filtering.
\end{itemize}

\subsection{Metrics}

We report accuracy, macro-F1, macro precision and recall, calibration error, and expected calibration error. For operational characterization, we also report latency, throughput, and cost with and without caching.

\section{Results}

Table~\ref{tab:ranking} shows the main comparative results on the 51-claim test split. The full system does not outperform the strongest simple baselines on this benchmark.

\begin{table}[htbp]
\centering
\caption{Comparative ranking on the 51-claim test split.}
\label{tab:ranking}
\begin{tabular}{lrrrrr}
\toprule
System & Accuracy & 95\% CI & Avg. Confidence & Calibration Error & ECE \\
\midrule
Random & 0.373 & [0.240, 0.505] & 47.7 & 0.104 & 0.134 \\
Majority & 0.353 & [0.222, 0.484] & 50.0 & 0.147 & 0.147 \\
Length & 0.314 & [0.186, 0.441] & 49.0 & 0.176 & 0.176 \\
Keyword & 0.294 & [0.169, 0.419] & 32.5 & 0.031 & 0.165 \\
Sentiment & 0.294 & [0.169, 0.419] & 40.0 & 0.106 & 0.106 \\
Ablate semantic reranking & 0.235 & [0.119, 0.352] & 40.6 & 0.170 & 0.212 \\
Full system & 0.216 & [0.103, 0.329] & 47.9 & 0.263 & 0.295 \\
Ablate debate & 0.216 & [0.103, 0.329] & 47.9 & 0.263 & 0.295 \\
Ablate quality filtering & 0.196 & [0.087, 0.305] & 39.2 & 0.196 & 0.196 \\
Ablate credibility & 0.137 & [0.043, 0.232] & 35.1 & 0.214 & 0.214 \\
\bottomrule
\end{tabular}
\end{table}

The ablation results are more informative than the headline ranking. Removing credibility scoring reduces accuracy from 21.6\% to 13.7\%, suggesting that the rubric contributes useful signal even though it is not sufficient to dominate all baselines. Removing semantic reranking increases accuracy slightly in this split, which implies that the reranking module is unstable under the current benchmark distribution. Debate has a neutral effect here.

\begin{table}[htbp]
\centering
\caption{Ablation summary.}
\label{tab:ablation}
\begin{tabular}{lrrr}
\toprule
Variant & Accuracy & Macro-F1 & Delta vs. full \\
\midrule
Full system & 0.216 & 0.212 & -- \\
Ablate credibility & 0.137 & 0.102 & -0.078 accuracy \\
Ablate semantic reranking & 0.235 & 0.225 & +0.020 accuracy \\
Ablate debate & 0.216 & 0.212 & 0.000 accuracy \\
Ablate quality filtering & 0.196 & 0.177 & -0.020 accuracy \\
\bottomrule
\end{tabular}
\end{table}

Operationally, the system remains attractive. Baseline latency is 8.20 seconds per claim, while debate mode raises average latency to 72.00 seconds, a ratio of 8.78\,$\times$. The corresponding throughput drops from 439.02 claims/hour to 50.00 claims/hour. These numbers make clear that debate should be used selectively rather than by default.

\begin{table}[htbp]
\centering
\caption{Operational metrics.}
\label{tab:ops}
\begin{tabular}{lr}
\toprule
Metric & Value \\
\midrule
Baseline latency & 8.20 s \\
Debate latency & 72.00 s \\
Latency ratio & 8.78$\times$ \\
Baseline throughput & 439.02 claims/hour \\
Debate throughput & 50.00 claims/hour \\
Monthly cost without cache & USD 77.00 \\
Monthly cost with cache & USD 22.00 \\
Monthly savings & 71.4\% \\
\bottomrule
\end{tabular}
\end{table}

\section{Discussion}

The results support a careful interpretation. The project is not yet a state-of-the-art fact classifier on the current split, but it is a well-engineered system with measurable, explainable behavior. That is an important distinction for a journal submission. The paper can credibly claim transparency, reproducibility, and cost-aware deployment, while treating broader performance claims as provisional.

The benchmark itself is currently the main bottleneck. The test split is small, contains synthetic patterns, and is not strong enough to support definitive claims about superiority. This is why the manuscript should emphasize system design and operational realism rather than benchmark dominance.

The credibility rubric is the clearest positive signal in the ablation study. It does not solve the verification problem alone, but it improves the overall pipeline. Debate mode remains valuable as an optional inspection tool, not as the default inference path.

\section{Limitations}

The main limitations are straightforward. First, the benchmark is limited in size and realism. Second, the system depends on search ranking and therefore inherits selection bias toward mainstream indexed sources. Third, calibration is imperfect, so confidence scores should be interpreted as approximate rather than probabilistic guarantees. Fourth, debate arbitration can increase latency substantially and should not be enabled indiscriminately.

These limitations are not incidental; they define the honest boundary of the current contribution.

\section{Conclusion}

Fact Validator demonstrates a practical path toward transparent, source-aware fact checking. Its key contribution is not raw benchmark dominance but an integrated architecture that combines evidence retrieval, auditable credibility scoring, optional debate, caching, and production safeguards. On the current benchmark, the system is best framed as a defensible system paper with provisional quality results and strong operational evidence. Future work should expand the benchmark with human-annotated claims, improve calibration, and evaluate selective debate policies on larger, cleaner datasets.

\begin{thebibliography}{99}

\bibitem{galitsky2023truthometer}
Galitsky, B. A. (2023).
Truth-O-Meter: Collaborating with LLM in Fighting its Hallucinations.
\textit{Preprints}, 202307.1723.
Available at: \url{https://www.preprints.org/manuscript/202307.1723}.

\bibitem{grave2020factool}
Grave, E., Bisk, Y., Alon, U., \& Holtzman, A. (2023).
FacTool: Factuality Detection in Generative AI via Tool-Augmented Framework.
\textit{ACL 2023 Findings}.

\bibitem{kubota2023zero}
Kubota, S., Kajiwara, Y., \& Onishi, T. (2023).
Zero-shot Faithful Factual Error Correction.
\textit{Proceedings of ACL 2023}.

\bibitem{thorne2018fever}
Thorne, J., Vlachos, A., Christodoulopoulos, C., \& Mittal, A. (2018).
FEVER: A Large-scale Dataset for Fact Extraction and Verification.
In \textit{Proceedings of NAACL-HLT}, pp. 809--819.

\bibitem{ji2023hallucination}
Ji, Z., Lee, N., Frieske, R., et al. (2023).
Survey of Hallucination in Natural Language Generation.
\textit{ACM Computing Surveys}, 55(12), 1--38.

\bibitem{wen2024abstention}
Wen, B., Yao, J., Feng, S., et al. (2024).
Know Your Limits: A Survey of Abstention in Large Language Models.
\textit{arXiv preprint arXiv:2407.18418}.

\end{thebibliography}

\end{document}